\section{A smaller cover from continuously varying placements}
\label{sec:adaptive}

The prism for each height may have its own orientation and translation.
This freedom gives the stronger endpoint
\[
                         V_3<0.052942305538 .
\]
The certificate below specifies a convex polytope and a rigid placement
for every height. Its volume check uses integers and fractions.

Write $(p,3h)=((1-t^2)/(1+t^2),2t/(1+t^2))$, with $0\le t\le1$,
and let $C$ be the quadrilateral from Section~\ref{sec:prism}.
The spatial prism at parameter $t$ has the eight corners
$(pa,pb,h\epsilon)$ for $(a,b)$ a vertex of $C$ and
$\epsilon\in\{0,1\}$. The alternating-plane argument already proves
that these prisms suffice for every unit curve.

The numerical data consist of 32 ordered rows in
\texttt{upper/poses.json}. Divide every entry by $D=10^8$.
A row lists $t$, four quaternion coordinates $(x,y,z,w)$ with scalar
part $w$ last, and a translation $g\in\mathbb R^3$. The parameters
increase from zero to one with one repeated value. The two rows at
that value give alternative placements at the change of orientation.
For each pair of rows with distinct consecutive parameters,
interpolate $t$, the quaternion $q$, and $g$ affinely in a common
variable $s\in[0,1]$. The adjacent quaternion dot products are
positive, so $q(s)$ never vanishes. For each corner use the map
\[
 (pa,pb,h\epsilon)\longmapsto
 R_q(pa,pb,h\epsilon)+g,\qquad
 R_q(v)=qv\overline q/\|q\|^2 .
\]
Quaternion conjugation here acts on vectors identified with pure
quaternions. The same proper rigid motion applies to all eight corners.

To enclose the moving corners, divide each interpolation interval into
three equal subintervals. In each one, parameterized by $u\in[0,1]$,
a moving corner has the form $N(u)/W(u)$, with
\[
 W=(1+t^2)\|q\|^2,\qquad
 N=q\big((1-t^2)a,(1-t^2)b,2t\epsilon/3\big)\overline q+Wg .
\]
All four homogeneous coordinates $(N,W)$ have degree at most five.
Expand them in the Bernstein basis
$B_j(u)=\binom5j u^j(1-u)^{5-j}$ with coefficients $(N_j,W_j)$,
$0\le j\le5$. The verifier evaluates at $u=0,1/5,\ldots,1$
and solves the interpolation equations in rational arithmetic.
It checks $W_j>0$ in every case. The identity
\[
 \frac{N(u)}{W(u)}
   =\sum_{j=0}^5\frac{B_j(u)W_j}{W(u)}\frac{N_j}{W_j}
\]
has nonnegative coefficients summing to one. The six Euclidean control
points therefore enclose the entire corner trajectory in their convex
hull, including every height between samples.

Let $P$ be the hull of the control points rounded coordinatewise to the
nearest multiple of $1/D$. It has 3,606 distinct generators. Set
$c=(0.17,0.13,0.07)$ and $r=1/40$. The facet check verifies
$c+[-r,r]^3\subseteq P$. With $e=1/(2D)$, convexity gives
\[
 P+[-e,e]^3\ \subseteq\ c+(1+e/r)(P-c).
\]
Indeed $[-e,e]^3\subseteq(e/r)(P-c)$, and the Minkowski sum of
positive homothetic copies of $P-c$ adds their scale factors. The
enlarged polytope on the right thus contains every original control
point, hence a rigid copy of every prism.

The program \texttt{upper/check\_adaptive.py} reconstructs and checks
all 2,955 facets from triples listed in
\texttt{upper/adaptive\_facets.json}. For each triple it checks the
integer supporting inequality against every generator, recovers the
entire polygon on the plane, and checks the eight cube corners by
evaluating the supporting normal on the cube. It verifies that every
boundary edge in the resulting list has two incident facets.
These are whole supporting facets, since each contains a noncollinear defining triple. The point set is full dimensional: a supporting triple has a point strictly on its inner side. In a three-polytope each edge belongs to exactly two facets. The graph joining facets across edges is connected, as one sees by projecting the boundary radially onto a sphere and joining facet interiors by a path avoiding vertices and crossing edges transversely. The verified nonempty family is closed under crossing edges and therefore contains all facets. Triangulating the facets
and summing the absolute tetrahedral determinants about $c$ gives
$|P|$. The final exact comparisons are
\[
 0.052942305537<(1+e/r)^3|P|<0.052942305538 .
\]

The file \texttt{upper/cover.obj} gives the vertices and outward faces of the
enlarged polytope. Its coordinates are terminating decimals: for a grid point
$p$ and cube center $c$ written in integer grid units, the dilated physical
coordinate is $(5{,}000{,}001p-c)/(5{,}000{,}000\cdot10^8)$.
Thus the mesh describes the certified body exactly.
